% Originally: new section (no predecessor in the week-based skeleton).
% Generator: Claude Opus 5 (High) --- \cref{sec:geodesics} opens by promising that "the
% essentially same calculation can be used to find the shortest path between points on a sphere
% or any other submanifold", and then never carries it out: the chapter ends with the flat case.
% This section closes that gap, and it is also the reason the chapter was moved to sit after
% \cref{ch:submanifolds} rather than in front of the vector-calculus material.

\section{Geodesics on a submanifold}
\label{sec:geodesics_on_submanifolds}

What survives of \qt{shortest means straight} when the curve is not free to leave a surface? On
$\mathbb{R}^n$ the answer was $c'' \equiv 0$, and on a sphere that is hopeless: no curve with
zero acceleration stays on a sphere at all. Something weaker has to replace it, and the
computation itself says what.

The only place the ambient space entered the proof of \cref{prop:geodesic_straight_line} was in
the choice of perturbations $g$. There, $g$ ranged over \emph{all} smooth fields vanishing at the
endpoints, and that is exactly what forced $c'' \equiv 0$. On a submanifold the competitors must
stay on $M$, so $g$ may only push along $M$, and the conclusion weakens in step: $c''$ has to be
orthogonal to the directions $g$ is allowed to take, and to nothing else.

\begin{proposition}[The geodesic equation]
\label{prop:geodesic_equation_submanifold}
Let $M \subseteq \mathbb{R}^n$ be a $C^2$ submanifold and let $c \in C^2([0,1],\mathbb{R}^n)$ with
$c(t) \in M$ for all $t$ and constant speed $\lvert c'(t)\rvert \equiv \lambda > 0$. If $c$
minimises length among $C^1$ curves in $M$ with the same endpoints, then
\[c''(t) \ \perp\ T_{c(t)}M \qquad \text{for every } t \in [0,1],\]
that is, $c''(t) \in N_{c(t)}M$ in the sense of \cref{def:normal_space}.
\end{proposition}

\begin{proof}[Proof sketch]
Let $g \in C^1$ vanish at the endpoints and be \newterm{tangential} along $c$, meaning
$g(t) \in T_{c(t)}M$ for every $t$. Take any family $c_\varepsilon$ of $C^1$ curves \emph{in $M$}
with $c_0 = c$, with the endpoints held fixed, and with
$\partial_\varepsilon\big|_{\varepsilon=0}c_\varepsilon = g$. Such a family exists: read $M$ near
$c(t)$ through a local parametrization (\cref{thm:local_parametrization}), push $g$ back to the
parameter domain, move there, and push forward again. This is the one step taken on trust here,
and it is the only place the submanifold structure is used.

From there the computation of \cref{prop:geodesic_straight_line} runs verbatim, since it never
asked where $g$ came from, and yields
\[0 = \left.\frac{d}{d\varepsilon}\right|_{\varepsilon=0}L(c_\varepsilon)
    = -\frac1\lambda\int_0^1 \big\langle c''(t),\,g(t)\big\rangle\,dt.\]
The fundamental lemma of the calculus of variations now applies to tangential $g$ only, so it
gives that the tangential \emph{part} of $c''(t)$ vanishes for every $t$, rather than $c''$
itself. That is precisely $c''(t) \perp T_{c(t)}M$.
\end{proof}

\begin{remark}[Reading the equation]
The statement has a mechanical reading worth carrying: a geodesic is a curve that accelerates
\emph{only as much as the surface forces it to}. Any acceleration along $M$ would be wasted
motion that a shorter competitor avoids, so all that is left points out of $M$, and the surface
supplies exactly that much to keep the curve on itself. Setting $M = \mathbb{R}^n$ recovers
\cref{prop:geodesic_straight_line}, since then $T_{c(t)}M = \mathbb{R}^n$ and orthogonality to
everything means $c'' \equiv 0$.

The converse fails, and for the same reason it failed in \cref{prop:extremum_implies_critical}:
the equation is a first-order criticality condition, so it detects candidates. On the sphere the
long way round a great circle satisfies it and is not shortest.
\end{remark}

% Generator: Claude Opus 5 (High)
\begin{aiexample}[Great circles, and why a circle of latitude is not a geodesic]
\label{ex:great_circles_geodesics}
Take $M := S^2 = \{x \in \mathbb{R}^3 : \lvert x\rvert^2 = 1\}$. By
\cref{prop:gradient_orthogonality_level_sets} applied to $g(x) := \lvert x\rvert^2$, the normal
space at $p \in S^2$ is $N_pS^2 = \operatorname{span}\{p\}$, so
\cref{prop:geodesic_equation_submanifold} reads: \textbf{$c$ is a geodesic candidate exactly when
$c''(t)$ is a multiple of $c(t)$.}

\textbf{A great circle satisfies it.} Fix orthonormal $p,v \in \mathbb{R}^3$ and put
\[c(t) := \cos(\lambda t)\,p + \sin(\lambda t)\,v.\]
Then $\lvert c(t)\rvert = 1$, so $c$ runs in $S^2$, and $\lvert c'(t)\rvert \equiv \lambda$, so
the constant-speed hypothesis holds. Differentiating twice,
\[c''(t) = -\lambda^2\big(\cos(\lambda t)\,p + \sin(\lambda t)\,v\big) = -\lambda^2\,c(t),\]
which is a multiple of $c(t)$. The condition holds at every $t$.

\textbf{A circle of latitude does not.} Fix a height $h$ with $0 < \lvert h\rvert < 1$, put
$r := \sqrt{1-h^2}$, and run around the circle at that height at constant speed:
\[c(t) := \big(r\cos(\omega t),\ r\sin(\omega t),\ h\big), \qquad \omega \neq 0.\]
Again $\lvert c(t)\rvert^2 = r^2+h^2 = 1$ and $\lvert c'(t)\rvert \equiv r\lvert\omega\rvert$ is
constant. But
\[c''(t) = -\omega^2\big(r\cos(\omega t),\ r\sin(\omega t),\ 0\big),\]
whose third coordinate is $0$ while the third coordinate of $c(t)$ is $h \neq 0$, and whose first
two coordinates are not both zero. So, $c''(t)$ is not a multiple of $c(t)$, and no circle of
latitude other than the equator ($h=0$, where the two agree and the curve is a great circle) is a
geodesic.

This is the flight-path fact in the form the equation gives it: flying due west along a parallel
is not the shortest route between two points on it, which is why long-haul routes bend towards
the pole.
\end{aiexample}

% Generator: Claude Opus 5 (High) --- FIG-CH18-01. The chapter had one figure in 300 lines, and
% none for the equation it exists to prove.
%
% A first version of this figure was drawn as ONE sphere carrying both curves, and rendering it
% killed it. Three faults, none visible in the source: the normal component of c'' lies along
% the radius by definition, so its arrow was drawn exactly on top of the radius line and the two
% collapsed into one unreadable stroke; the "tangential part" and "N_p S^2" labels overlapped;
% and the great circle was drawn with no vector on it at all, so it asserted "c'' is radial"
% while showing nothing. Two panels, and no radius line, fix all three.
%
% GEOMETRY, checked, so the drawing can be audited without redoing it. Both panels are drawn in
% the plane y = 0, so screen coordinates ARE the (x,z) coordinates of R^3 and the vectors need
% no projection; only the circles that leave that plane are foreshortened into ellipses.
% Sphere radius 1.6 in both.
%
% LEFT (great circle). The equator is itself a great circle, so it is the one drawn. Marked
% point p := (-1.6, 0, 0), on the silhouette. For c(t) = R(cos t, sin t, 0) one has c'' = -c,
% i.e. it points straight at the centre: purely normal, nothing tangential left over.
%
% RIGHT (circle of latitude). Height h = 0.96, hence radius r = sqrt(2.56 - 0.9216) = 1.28;
% ellipse centred (0,0.96) with semi-axes 1.28 and 0.40. Marked point p := (-1.28, 0, 0.96),
% which satisfies |p| = sqrt(1.6384 + 0.9216) = 1.6, so it is on the sphere, and it is both the
% leftmost point of the latitude ellipse and on the silhouette. That is what lets it be drawn
% honestly. Unit normal n = p/1.6 = (-0.8, 0, 0.6). The acceleration points at the axis, along
% +x; take c'' = (1.1, 0, 0). Then
%   normal part      = <c'',n> n = (-0.88)(-0.8,0,0.6) = (0.704, 0, -0.528),
%   tangential part  = c'' - that = (0.396, 0, 0.528),
% and the check <(0.396,0,0.528), (-0.8,0,0.6)> = -0.3168 + 0.3168 = 0 confirms the second is
% genuinely tangent, while the two sum to (1.1,0,0) as they must. Tips, from p:
%   c''  -> (-0.18, 0.96);  normal -> (-0.576, 0.432);  tangential -> (-0.884, 1.488).
% The whole point of the figure is that the last of these is not p itself.
\begin{center}
\begin{tikzpicture}[>=Latex, scale=1.05]
  % ---- LEFT: a great circle. Its acceleration points at the centre. ----
  \begin{scope}
    \draw[gray!55] (0,0) circle (1.6);
    \draw[blue!70!black, thick] (0,0) ellipse (1.6 and 0.5);
    \fill (0,0) circle (1.1pt);
    \draw[red!75!black, thick, ->] (-1.6,0) -- (-0.32,0);
    \node[red!75!black, font=\scriptsize, anchor=south] at (-0.95,0.06) {$c''$};
    \fill[blue!70!black] (-1.6,0) circle (1.5pt);
    \node[blue!70!black, font=\scriptsize, anchor=east] at (-1.68,0) {$p$};
    \node[font=\scriptsize, anchor=north, align=center] at (0,-1.95)
      {great circle: $c''$ points at the centre,\\so it is purely normal};
  \end{scope}

  % ---- RIGHT: a circle of latitude. A tangential part survives. ----
  \begin{scope}[shift={(4.7,0)}]
    \draw[gray!55] (0,0) circle (1.6);
    \draw[gray!45, dashed] (0,0) ellipse (1.6 and 0.5);
    \draw[blue!70!black, thick] (0,0.96) ellipse (1.28 and 0.40);
    \fill (0,0) circle (1.1pt);

    % Guide line along the radius, i.e. along the normal direction, drawn thin and dotted so
    % that the normal-component arrow reads as lying on it rather than fighting with it.
    \draw[gray!45, dotted] (-1.28,0.96) -- (0.64,-0.48);

    % The decomposition. Parallelogram first, so the arrows sit on top of it.
    \draw[gray!55, dotted] (-0.576,0.432) -- (-0.18,0.96);
    \draw[gray!55, dotted] (-0.884,1.488) -- (-0.18,0.96);
    \draw[gray!50!black, dashed, ->] (-1.28,0.96) -- (-0.576,0.432);
    \node[gray!45!black, font=\scriptsize, anchor=north west] at (-0.62,0.40) {normal};
    \draw[OliveGreen, very thick, ->] (-1.28,0.96) -- (-0.884,1.488);
    \draw[red!75!black, thick, ->] (-1.28,0.96) -- (-0.18,0.96);
    \node[red!75!black, font=\scriptsize, anchor=west] at (-0.12,0.96) {$c''$};
    \node[OliveGreen, font=\scriptsize, anchor=south] at (-0.55,1.60) {tangential part};

    \fill[blue!70!black] (-1.28,0.96) circle (1.5pt);
    \node[blue!70!black, font=\scriptsize, anchor=east] at (-1.36,0.96) {$p$};
    \node[font=\scriptsize, anchor=north, align=center] at (0,-1.95)
      {circle of latitude: a tangential\\part survives, so it is no geodesic};
  \end{scope}
\end{tikzpicture}
\end{center}

% Generator: Claude Opus 5 (High)
\begin{aiexercise}[Geodesics on the cylinder]
\label{ex:ai_cylinder_geodesics}
Let $Z := \{(x,y,z) \in \mathbb{R}^3 : x^2+y^2 = 1\}$ be the unit cylinder, and for
$a, b \in \mathbb{R}$ not both zero consider the helix
\[c(t) := \big(\cos(at),\ \sin(at),\ bt\big).\]
\begin{enumerate}[label=\textbf{(\alph*)}]
  \item Show that $c$ runs in $Z$ at constant speed, and compute that speed.
  \item Determine $N_{c(t)}Z$, and verify that $c$ satisfies the geodesic equation of
    \cref{prop:geodesic_equation_submanifold} for every choice of $a$ and $b$.
  \item Which curves do the cases $a = 0$ and $b = 0$ describe?
\end{enumerate}
\end{aiexercise}
